% Cite Commanda
% \citet{}
% \citep{}
% \citet*{}
% \citep*{}
% \citeauthor{}
% \citeyear{}



\section{Notes on papers}


\subsection{\cite{fiordelmondo_towards_2025}}
``The need for documentation that goes beyond just publishing an article (which is still an important document) has been ad- vocated to promote critical discourse [3, 8]''

``documentation allows the critical discourse and the research field to grow and mature.''

[3] S. M. Astrid Bin. 2021. Discourse is critical: Towards a collaborative NIME
history. In Proceedings of the International Conference on New Interfaces for Musical Expression. Shanghai, China, Article 8. https://doi.org/10.21428/ 92fbeb44.ac5d43e1

[8] Filipe Calegario, João Tragtenberg, Christian Frisson, Eduardo Meneses, Joseph
Malloch, Vincent Cusson, and Marcelo M. Wanderley. 2021. Documentation and Replicability in the NIME Community. In Proceedings of the International Conference on New Interfaces for Musical Expression. Shanghai, China, Article 4. https://doi.org/10.21428/92fbeb44.dc50e34d

\subsection{papers citing \cite{mcpherson_nimehub_2016}}


\item Digital Musical Instruments as Research Products
\item Contextualizing Musical Organics: An Ad-hoc Organological Classification Approach
\item Rebuilding and Reinterpreting a Digital Musical Instrument - The Sponge

\cite{meneses_puara_2023} considers the performer, but what of the developer, or developer / performer?

\begin{quote}
    For maintaining instruments, we have, in addition to hardware replicability [3], a software layer, usually in the form of
firmware, to make the hardware work. Maintaining code in
a research-oriented environment where the creator is often a
student that will eventually move on to another project or
topic is difficult. The lack of team continuity is particularly
detrimental to instrument longevity.
\end{quote}


\begin{quote}
In addition, artists and projects supporting Free and OpenSource Software (FOSS) face another issue: FOSS usually
focuses on functionality but is not necessarily accessible to
non-tech-savvy users. I.e., high entry fee to use the tool,
hard to install/experiment, e.g., LivePose [9]; depends on additional knowledge such as a text-based music programming
language, e.g., SATIE [1].
\end{quote}


\cite{morreale_design_2017}

\begin{quote}
    “it is more sustainable, it allows derivatives
and offers inspiration to proliferate”.
\end{quote}



\cite{west_towards_2025}

\begin{quote}
    Recent discussions have highlighted the benefits of replication
[1, 7] and of engaging with old instruments more broadly [2]
towards the well-being of NIME practice and research. One major
theme in these discussions is the difficulty of documentation; it
is challenging to anticipate what information will be most useful
to later users and researchers who may wish to replicate an
instrument
\end{quote}

would that process have been facilitated by access to the cad files?


There is the pragmatic gap of what is explained, what is intended and what is understood, in particular by code.


\cite{ajin_rebuilding_2019}



\cite{calegario_documentation_2021}


\begin{quote}
    Replication is also essential in learning since replicating existing works allows us to understand complex developments' intricacies fully.
\end{quote}


\cite{burns_documents_2021}

\begin{quote}
    the open science community, which
seems to largely assume that more documentation of the parts of doing science produces a
stronger science alone
\end{quote}

I think this might be a misreading of open science, the intention in the first instance is the sharing of material \textit{that already exists},

\begin{quote}
    One important question is whether the web and internet
technologies that exist,48 and the type of documents involved in disseminating explicit know-
how and know-that knowledge, are embedded within a culture that welcomes treating the
problem-finding dilemma with the attention it deserves.
\end{quote}

\begin{quote}
    This reductionist way of explaining things by lieu of their physical components (and of
viewing the world, i.e., defining reality or what is real [the ontological]) was a hallmark of
positivism and has been common across all domains of knowledge and areas of practice. Several
years ago I was doing historical work at an institutional archive and reading annual library
reports from around the mid-twentieth century.37 One of the common problems that the head
librarian described in those reports concerned the university administration’s view of the library
as nothing more than a warehouse of books. As the head librarian, Ralph Parker, described it, the
administration at his institution ignored (or even failed to see) the complexities attached to
managing and using a library and, as a result, repeatedly failed to invest in the library and the
librarians who operated it (the practical implication of a reductionist viewpoint). For the
librarian, the library was more than a warehouse of books; in the process of acquiring,
describing, managing, shelving, circulating, and using the works in the library, and the library
itself, something greater emerged or came into existence that was beyond a basic warehouse.
And the thing that emerged was as real as any of its constitutive parts (e.g., books and shelves),
even if it could not be reduced to those parts. Thus, even though the administration would agree
that the library was a real place, it was only real to them as the “cobblestones” of the street are
real—as opposed to whole streets with their social worlds and other components (the whole is
reduced to a part that then becomes its perceived entire ontological substance; perhaps we can
think of this as synecdochic fallacy38)
\end{quote}


The intention is not to be reductionist, to look at components


\begin{quote}
     If so, is it tacit
knowledge rather than just implicit in some way; is it the case that there will always be some
kind of knowledge that will forever resist documentation; and whether, among other things, is it
worth documenting, or attempting to document, certain kinds of tacit or implicit knowledge,
because, for example, the cost is high in some important way or the cognitive load of too much
documentation is burdensome
\end{quote}

Neither is the goal to place a burden of documentation.


\cite{perkel_democratic_2016} cited in \cite{national_open_2018}


\cite{burns_tacit_2021}

\begin{quote}
    "Open methods have the same effect [ as open data ] and also facilitate
progress in reuse, adaptation, and extension for new research" (Nosek et al.,
p. 624, 2012).
"Authors cannot identify and report [ in publications ] every detail that may be
important in a method, but many more parts of the methodology can be
shared outside of the report itself" (Nosek et al., p. 625, 2012).
\end{quote}

Nosek, B. A., Spies, J. R., & Motyl, M. (2012). Scientific Utopia: II. Restructuring Incentives and Practices to Promote Truth
Over Publishability. Perspectives on Psychological Science, 7(6), 615–631. https://doi.org/10.1177/1745691612459058

\begin{quote}
    Technologies in place to support open science work flows
● Electronic Lab Notebooks
○ LabArchives: https://www.labarchives.com/
○ OSF: https://osf.io/
○ Project Jupyter: https://jupyter.org/
● Data Repositories: Recommended Data Repositories from Nature.com
● Manuscript Repositories: arXiv.org and others listed on Wikipedia
● Code Repositories: GitHub, GitLab, and more
● Social Media (general and specifi
\end{quote}





% Dependencies!

% Missing but extant repositories,
% for which burden of locating them is placed on the reader, ambiguity, no longer findable and accessible
% Impact on longevity

% The burden to get it right is placed I. The author

% A suggested methodology through the case study of a DMI.

% \subsubsection{Missing Repos [Notes]}
% \url{https://gitea.offig.com/lfsadmin/Cicadas.git} from \textit{A Synthetic Cicada Soundscape Controlled by Breath} does \textit{not} count as it is behind authentication.
% \url{https://github.com/ijc8/glitchgate} as it is unreachable (and likely just still marked private)
% \textit{Synthesizing Music with Logic Gate Networks}
% \begin{enumerate}
% \item Live Improvisation with Fine-Tuned Generative AI: A Musical Metacreation Approach \url{https://github.com/MAz-Codes/sao-silence-remover} is only part of the process, but non-zero
% \end{enumerate}
% sonicolour with a link to a google site to a link to a git repository that contains a zip file \textit{does} count
% \url{https://github.com/Samson026/LibPdIntegrationForAndroid} from \textit{Luna: An AR Musical Instrument on the Meta Quest 2} does not count as it it a constituent part, but not the whole
% \textit{Roger Dannenberg and Ari Liloia. 2022. Global Drum Circle (Private Github Repository).} is that allowed
% Exploiting Latency In The Design Of A Networked Music Performance System For Percussive Collective Improvisation
% This shouldn't be allowed:
% [15] Roger Dannenberg and Ari Liloia. 2022. Global Drum Circle (Private Github Repository).

% but something like \url{tinytouchinstruments.com} is incidentally open given the nature of Javascript so does not count.

% \url{https://github.com/mayacaren/ViVo} missing dependencies, for instance, thats fine, there is at least something 


% Similar questions have been asked previously. \cite{calegario_documentation_2021} outlines what should be included and \citet{fiordelmondo_longevity_2024}  provides a template based up on it to reduce the cognitive labour to carry it out.
% Both informed the approach taken to searching through material.

% \citet{calegario_documentation_2021} and \citet{fiordelmondo_towards_2025} exploring the usage and dissemination of open source projects at NIME. 

% Only 2025 was considered as it is the most up to date digital research practice can be.


% \cite{calegario_documentation_2021} outlines what should be included and \citet{fiordelmondo_longevity_2024}  provides a template based up on it to reduce the cognitive labour to carry it out.

% An empty repository would be disqualidfied.
% Source material defined as the material pertinent to the paper without, in a version control format, digital archive or ideally both.
% To begin answering this question, NIME papers from the last 10 years of NIME proceedings were downloaded.
% All 96 papers from NIME 2025 can be downloaded locally with a zshell instruction
% \begin{verbatim}
% for i in {1..96}; do
%   curl "https://nime.org/proceedings/2025/nime2025_${i}.pdf" -O
% done
% \end{verbatim}


% The process takes about three hours to complete


